\section{Query Dasar: SELECT, INSERT, UPDATE, DELETE}

Setelah koneksi MySQLi terbentuk, operasi data pada database dilakukan melalui \textbf{query SQL}. Empat operasi dasar yang dikenal sebagai \textbf{CRUD} (\textit{Create, Read, Update, Delete}) dipetakan ke perintah SQL: \code{INSERT} (create), \code{SELECT} (read), \code{UPDATE} (update), dan \code{DELETE} (delete). Query dijalankan menggunakan metode \code{\$conn->query(\$sql)} untuk query tanpa parameter, atau menggunakan \textit{prepared statement} untuk query yang melibatkan input pengguna --- ini penting untuk mencegah \textbf{SQL injection} \cite{phpManual}.

\code{SELECT} digunakan untuk membaca data dari tabel. Hasilnya berupa \textit{result set} yang dapat diiterasi menggunakan metode \code{fetch\_assoc()} (mengembalikan satu baris sebagai array asosiatif), \code{fetch\_all(MYSQLI\_ASSOC)} (semua baris sekaligus), atau \code{fetch\_object()}. Properti \code{\$result->num\_rows} menunjukkan jumlah baris. Untuk operasi \code{INSERT}, \code{UPDATE}, dan \code{DELETE}, metode \code{query()} mengembalikan boolean, dan \code{\$conn->affected\_rows} menunjukkan jumlah baris yang terpengaruh. Setelah \code{INSERT}, \code{\$conn->insert\_id} berisi ID yang baru digenerate (auto-increment) \cite{mysqlDocs}.

\begin{webcode}[language=PHP, caption={CRUD dengan prepared statement (MySQLi OOP)}]
<?php
$conn = new mysqli('localhost', 'root', '', 'belajar_php');
$conn->set_charset("utf8mb4");

// CREATE - INSERT dengan prepared statement
$stmt = $conn->prepare("INSERT INTO tugas (judul, deskripsi) VALUES (?, ?)");
$stmt->bind_param("ss", $judul, $deskripsi);
$judul     = "Belajar PHP OOP";
$deskripsi = "Memahami class, object, dan inheritance";
$stmt->execute();
echo "Data ditambahkan, ID: " . $stmt->insert_id . "\n";
$stmt->close();

// READ - SELECT
$result = $conn->query("SELECT id, judul, deskripsi FROM tugas ORDER BY id");
while ($row = $result->fetch_assoc()) {
  echo "#{$row['id']} {$row['judul']}: {$row['deskripsi']}\n";
}
$result->free();

// UPDATE dengan prepared statement
$stmt = $conn->prepare("UPDATE tugas SET judul = ? WHERE id = ?");
$stmt->bind_param("si", $judulBaru, $id);
$judulBaru = "Belajar PHP OOP Lanjutan";
$id = 1;
$stmt->execute();
echo "Baris diupdate: " . $stmt->affected_rows . "\n";
$stmt->close();

// DELETE dengan prepared statement
$stmt = $conn->prepare("DELETE FROM tugas WHERE id = ?");
$stmt->bind_param("i", $id);
$id = 1;
$stmt->execute();
echo "Baris dihapus: " . $stmt->affected_rows . "\n";
$stmt->close();

$conn->close();
?>
\end{webcode}

\begin{table}[ht]
\centering
\begin{tabular}{|P{2cm}|P{3cm}|P{3cm}|P{4cm}|}
\hline
\textbf{Operasi} & \textbf{SQL} & \textbf{Return} & \textbf{Info tambahan} \\
\hline
Create & \code{INSERT} & Boolean & \code{\$conn->insert\_id} \\
\hline
Read & \code{SELECT} & Result set & \code{fetch\_assoc()}, \code{num\_rows} \\
\hline
Update & \code{UPDATE} & Boolean & \code{\$conn->affected\_rows} \\
\hline
Delete & \code{DELETE} & Boolean & \code{\$conn->affected\_rows} \\
\hline
\end{tabular}
\caption{Ringkasan operasi CRUD dengan MySQLi}
\end{table}

\textbf{Prepared statement} adalah mekanisme yang memisahkan query SQL dari data. Langkahnya: (1) \textit{prepare} template query dengan placeholder \code{?}; (2) \textit{bind} parameter ke placeholder menggunakan \code{bind\_param()} dengan indikator tipe (\code{"s"} = string, \code{"i"} = integer, \code{"d"} = double, \code{"b"} = blob); (3) \textit{execute} query. Prepared statement mencegah SQL injection karena data diperlakukan sebagai literal, bukan bagian dari perintah SQL. Selain keamanan, prepared statement juga memberikan keuntungan performa saat query yang sama dieksekusi berulang kali dengan data berbeda \cite{phpRightWay}.

\begin{catatan}
Jangan pernah menggabungkan string input langsung ke SQL seperti \code{"SELECT * FROM user WHERE id = \$id"}. Ini membuka celah SQL injection. Selalu gunakan prepared statement atau minimal \code{\$conn->real\_escape\_string()}, meskipun prepared statement jauh lebih aman dan direkomendasikan.
\end{catatan}

